\clearpage
\begin{summarybox}
\textbf{C090\quad 共焦椭圆与双曲线的焦角--离心率关系：快记／快懂／快用}

\textbf{快记（一条恒等式＋一条推论）：}
\[
 \boxed{\frac{\sin^2(\theta/2)}{e_1^2}+
 \frac{\cos^2(\theta/2)}{e_2^2}=1},\qquad
 \boxed{\frac1{e_1}+\frac1{e_2}\le\frac2{\sin\theta}}.
\]
其中 $0<e_1<1<e_2$，$0<\theta<\pi$。

\textbf{快懂（先看和差，不要先写坐标）：}
\[
 PF_1+PF_2=2a_E,\quad |PF_1-PF_2|=2a_H
 \quad\Longrightarrow\quad PF_1,PF_2=a_E\pm a_H.
\]
\textbf{为什么是 $a_E\pm a_H$？}两腰的\emph{平均长度}为 $a_E$，
离平均长度的\emph{偏差}为 $a_H$。余弦定理把两腰及夹角
$\theta$ 与底边 $2c$ 连起来；换半角后得到
\[
c^2=a_E^2\sin^2\frac\theta2+a_H^2\cos^2\frac\theta2.
\]
这是两种轴长平方的\emph{加权平均}，故 $a_H<c<a_E$。

\textbf{快用（考场 3 步）：}
\begin{enumerate}
\item 见到\emph{共焦＋公共点＋焦角}，先画 $\triangle F_1PF_2$；
\item 用“和、差解两腰”，写出 $a_E\pm a_H$；
用余弦定理连接 $2c$ 和 $\theta$，再代入离心率；
\item 若要求 $1/e_1+1/e_2$，配对
$(\sin\frac\theta2/e_1,\cos\frac\theta2/e_2)$ 与
$(1/\sin\frac\theta2,1/\cos\frac\theta2)$ 用柯西。
\end{enumerate}

\textbf{等号核查：}
\[
 \boxed{e_1=\tan\frac\theta2,\quad e_2=\cot\frac\theta2,
 \quad0<\theta<\frac\pi2}.
\]
\textbf{固定角度的可行范围及边界：}
\[
\boxed{\sin\frac\theta2<e_1<1},\qquad e_2>1.
\]
左端只对应 $e_2\to+\infty$；右端对应 $e_1\to1^-,e_2\to1^+$；
两端都\emph{不可实际取到}。
\textbf{等号的几何样子：} $e_1e_2=1$ 时，公共点落在过某个焦点
且垂直焦点轴的直线上；\emph{直角在焦点，不在 $P$}。
$\theta\ge90^\circ$ 时 $1/e_1+1/e_2<2$（仅有上确界 $2$）。
\textbf{一句话记牢：}「椭圆给和，双曲线给差，余弦定理搭桥；
柯西得上界，等号须锐角。」
\end{summarybox}
